\documentclass[10pt,a4paper]{amsart}
\usepackage[margin=19mm]{geometry}
\usepackage{amsmath,amssymb,mathtools}
\usepackage[scr=rsfs,cal=euler]{mathalfa}
\usepackage{xfrac}
\usepackage[unicode,hidelinks,bookmarks=false]{hyperref}
\newcommand{\ie}{i.e.\ }
\newcommand{\cf}{cf.\ }
\newcommand{\resp}{respectively\ }
\newcommand{\Subsection}{\subsection}
\providecommand{\nobreakdash}{\nobreak\hbox{-}}
\setlength{\emergencystretch}{2em}
\multlinegap0pt
\usepackage{bm}
\usepackage{bbm}
\usepackage{dsfont}
\usepackage{xspace}
\usepackage{cancel}
\usepackage{mathtools}
\usepackage{amsthm}
\makeatletter
\@ifundefined{theorem}{\newtheorem{theorem}{Theorem}}{}
\@ifundefined{corollary}{\newtheorem{corollary}[theorem]{Corollary}}{}
\@ifundefined{lemma}{\newtheorem{lemma}[theorem]{Lemma}}{}
\makeatother
\newcommand{\J}{\mathcal{J}}
\newcommand{\N}{\mathbb{N}}
\newcommand{\R}{\mathbb{R}}
\title[A Cubic Regularization Method for Multiobjective Optimizatio]{A Cubic Regularization Method for Multiobjective Optimization}
\author{Douglas S. Gon\c{c}alves, Max L. N. Gon\c{c}alves, Jefferson G. Melo}
\date{Source: arXiv:2506.08181v1 (CC BY 4.0)}
\begin{document}
\maketitle

\noindent\emph{Source and licence.} A Cubic Regularization Method for Multiobjective Optimization. Version arXiv:2506.08181v1 (CC BY 4.0), distributed under the Creative Commons Attribution 4.0 International licence, as recorded in the arXiv licence metadata for this exact version and shown on the abstract page https://arxiv.org/abs/2506.08181v1. Full rights notice: licenses/arxiv-2506.08181-CC-BY-4.0.txt.\par\smallskip

\noindent\textit{Editorial scope.} Counted unit: the Theorem whose source label is \texttt{theorem:superlinearconvergence} in the article (source0.tex lines 853--865, source label \texttt{theorem:superlinearconvergence}), delivered together with the article's own complete original proof of it (source0.tex lines 866--921, one \texttt{proof} environment of 56 lines). No mathematics is written, repaired or re-derived: statement and proof are the article's own text line for line. Required context retained uncounted: vectorproblem (lines 240--242); definition-approxGrad-Hessian (lines 386--388); lem3:basic-gt+1<Deltax (lines 454--460); ineq:ApproxGrad-Hess (lines 502--505); eq:subprob-Alg1 (lines 516--519); def:sigmaMax (lines 564--566); cor:conv (lines 719--726); ineq:ApproxGrad-Hess12 (lines 763--766); def:gt-Deltaxt (lines 768--770); thm:3.3 (lines 780--796); eq:Deltaxt-lambdanablaFxt (lines 787--789); ineq:contraction gt (lines 836--838). Every definition, display and label of those lines is retained unchanged. Omitted from this package and not redistributed: 1--239, 243--385, 389--453, 461--501, 506--515, 520--563, 567--718, 727--762, 767--767, 771--779, 790--835, 839--852, 922--1622. Nothing inside a retained excerpt is omitted or altered: every retained body is delivered exactly as the article's lines are written, so no omission receipt of any kind accompanies this package. Citations: the retained excerpts carry no citation command at all, so no bibliography entry of the article is transcribed and no reference is redistributed. Reference adaptation: none was needed and none was made; every reference inside the delivered unit resolves inside it, so no unresolved reference or citation is produced. Printed slips kept literal: no printed word, symbol, hypothesis or number of the counted statement or of its proof was altered. Layout and class: the article's own class is \texttt{siamart250211} with packages including lipsum, amssymb, amsfonts, amstext, amsmath, graphicx, geometry, epstopdf, algorithmic, cite, url, appendix, multirow, multicol; the wrapper is the lane's pinned \texttt{amsart} editorial wrapper at 10pt on A4, which restates the theorem-like environments the retained text needs and adds the article's own macro definitions that the retained text needs (\texttt{\textbackslash J}, \texttt{\textbackslash N}, \texttt{\textbackslash R}), with extra wrapper packages (bm, bbm, dsfont, xspace, cancel, mathtools). The delivered numbering is the unit's own. Assets: no retained span contains a figure environment, a graphics-inclusion command, a PDF-page inclusion, a table or a TikZ picture; no image file is copied into this package, the package declares no asset dependency and no image file is redistributed with it. Licence: version arXiv:2506.08181v1 of this article is distributed under the Creative Commons Attribution 4.0 International licence, as recorded in the arXiv licence metadata for this exact version and shown on the abstract page https://arxiv.org/abs/2506.08181v1. A full rights notice is delivered at licenses/arxiv-2506.08181-CC-BY-4.0.txt. Characters: every retained excerpt is pure ASCII, so the delivered engine, XeLaTeX with its default Latin Modern text font, reports no missing character.
\begin{equation} \label{vectorproblem}
\min_{x \in \R^n}    F(x),  \\
\end{equation}

\begin{equation}\label{definition-approxGrad-Hessian}
E^G_{j}(x):=\nabla F_j(x)-\nabla\bar F_j(x),  \quad E^H_{j}(x):=\nabla^{2}F_j(x)-\nabla^{2}\bar F_j(x), \quad \forall j\in \J.
\end{equation}

\begin{lemma}
\label{lem3:basic-gt+1<Deltax}
Let $x\in \R^n$ be given and consider $E^G_{j}(x)$ and  $E^H_{j}(x)$ as in  \eqref{definition-approxGrad-Hessian}. Assume that {\bf (A1)} holds and let a triple $(x^{+},\lambda^{+},\sigma) \in \R^n\times \R^m_+\times\R_{++} $ be such that $\sum_{j=1}^{m} \lambda^+_j=1$ and $\left\|\sum_{j \in \J}\lambda^+_j\nabla M_{x,\sigma}^j(x^{+})\right\|\leq \theta\|x^+- x\|^2$, for some $\theta\geq 0$. Then, we have   
\begin{equation*}
\left\|\sum_{j \in \J}\lambda^+_j\nabla F_j(x^{+})\right\|\leq \frac{L+\sigma+2\theta}{2}\|x^{+}-x\|^2+\max_j\left(\|E_j^G(x)\|+ {\|E_j^H(x)\|\|x^{+}-x\|}\right). 
\end{equation*} 
\end{lemma}

\begin{equation}
\|E^G_j(x_t)\|\leq\frac{\bar \kappa_{G}}{\alpha^{i-1}}\|x_t-x_{t-1}\|^2,  \quad \|E^H_j(x_t)\|\leq\frac{\bar \kappa_{H}}{\alpha^{i-1}}\|x_t-x_{t-1}\|.
\label{ineq:ApproxGrad-Hess}
\end{equation}

\begin{equation}
M_{x_{t},\alpha^{i}\sigma_{t}}(x^{+}_{t,i})\leq 0, \quad \sum_{j=1}^{m} (\lambda^+_{t,i})_j=1, \quad \left \|\sum_{j=1}^{m} (\lambda^+_{t,i})_j\nabla M_{x_t,\alpha^i\sigma_t}^j(x_{t,i}^+)\right\|\leq \theta  \|x_{t,i}^+-x_t\|^2,
\label{eq:subprob-Alg1}
\end{equation}

\begin{equation}\label{def:sigmaMax}
\sigma_{1}\leq\sigma_{t}\leq \sigma_1+2[L+3\alpha(2\bar \kappa_{G}+\bar \kappa_{H})]:=\sigma_{max},
\end{equation}

\begin{corollary}\label{cor:conv}
Assume that {\bf (A1)} and {\bf (A2)} hold and let $\{x_{t}\}$ be generated by Algorithm~1. Then, 
\begin{equation}
\lim_{t\to+\infty} \|x_t-x_{t-1}\|=0, \qquad \lim_{t\to+\infty}\left \|\sum_{j\in\J}(\lambda_{t})_j\nabla F_j(x_{t})\right\|=0.
\label{eq:3.16}
\end{equation}
As a consequence, all limit points of $\{x_t\}$, if any, are weak Pareto critical for \eqref{vectorproblem}.
\end{corollary}

\begin{equation}
 \|E^H_j(x_t)\|\leq\frac{\bar \kappa_{H}}{\alpha^{i-1}}\min\left\{ \|\Delta x_{t}\|, \zeta\left\|g_t\right\| \right\}.
\label{ineq:ApproxGrad-Hess12}
\end{equation}

\begin{equation}\label{def:gt-Deltaxt}
 g_t:=   \sum_{j\in \J} (\lambda_{t})_j\nabla F_j(x_{t}), \qquad \Delta x_t:=x_t-x_{t-1}.
\end{equation}

\begin{theorem} \label{thm:3.3}
Assume that  {\bf (A1)-(A4)} hold and let $\left\{x_{t}\right\}$ be generated by Algorithm 1. Then, there exists $t_1\in \N$ such that, for all $\lambda \in \R^m_+$ such that $\sum_j\lambda_j=1$, there hold
% \begin{equation}
% \left\|\sum_{j\in \J}(\lambda_{t_0})_j\nabla F_j(x_{t_0})\right\|\leq\min\left\{\dfrac{\mu}{2\bar \kappa_{H}},\frac1{C}\right\},
% \label{eq:3.25}
% \end{equation}
% then the following inequalities hold
\begin{equation}\label{eq:Deltaxt-lambdanablaFxt} 
\|\Delta x_{t+1}\|\leq \frac{4}{\mu}\left\|\sum_{j\in \J}\lambda_j \nabla F_j(x_t)\right\|, \quad  \forall t\geq t_1,
\end{equation}
\begin{equation}
\left \|g_{t+1}\right\|\leq C\left\|g_t\right\|^{2},  \quad \forall t\geq t_1,
\label{eq:nablaFxt+1<=nablaFxt}
\end{equation}
where $C:= 8\left(L+\alpha\sigma_{max}+2\theta\right)/\mu^2+4\alpha\bar\kappa_H\zeta/\mu$.
Furthermore, the whole sequence $\{x_t\}$ converges to a local Pareto solution of \eqref{vectorproblem}.
\end{theorem}

\begin{equation} \label{ineq:contraction gt}
\|g_{t+1}\|\leq C\|g_t\|^2\leq (1/2)\|g_t\|, \quad \forall t\geq t_1.
\end{equation}

\begin{theorem}\label{theorem:superlinearconvergence}
Assume that {\bf (A1)--(A4)} hold and let $x^*$ be the limit point of $\{x_t\}$. Then,   there exist $C_1,C_2, C_3>0$ such that, for all $t\geq t_1$, there hold 
\begin{equation}\label{ineq:xi-xt<gt+1}
\|x_{t+1}-x^*\|\leq C_1\|g_{t+1}\| ,  
\end{equation} 
\begin{equation}\label{aux-eq: tilde gt}
  \left\|\tilde g_t\right\| \leq C_2 \min\{\|\Delta x_{t+1}\|, \|x_t-x^*\|\},
 \end{equation}
 \begin{equation}\label{ineq:xt+1-xbar leq gt xt-xbar}
  \left\|x_{t+1}-x^*\right\| \leq C_3 \|g_t\|\|x_t-x^*\|,
 \end{equation}
where $\tilde g_t:=\sum_{j\in \J} (\lambda_{t+1})_j\nabla F_j(x_t)$. As a consequence,  $\{x_t\}$ is superlinearly convergent  to $x^*$.
\end{theorem}

\begin{proof}
It follows from \eqref{eq:Deltaxt-lambdanablaFxt}  with $\lambda=\lambda_t$ and \eqref{ineq:contraction gt} that, for all $s\geq t+1$,
\begin{align*}
\|x_s-x_{t+1}\|&\leq \sum_{k=t+1}^s\|x_{k+1}-x_{k}\|=\sum_{k=t+1}^s\|\Delta x_{k+1}\|\nonumber\\
&\leq \frac{4}{\mu}\sum_{k=t+1}^s\|g_k\|=\frac{4}{\mu}\left[\|g_{t+1}\|+\|g_{t+2}\|+\cdots+ \|g_{s}\|\right]\nonumber\\
&\leq \frac{4}{\mu}\left[\|g_{t+1}\|+\frac{1}{2}\|g_{t+1}\|+\frac{1}{2^2}\|g_{t+1}\|+\cdots+ \frac{1}{2^{s-t-1}}\|g_{t+1}\|\right]\nonumber\\
&=\frac{4}{\mu}\left\{\sum_{l=0}^{s-t-1}\frac{1}{2^l}\|g_{t+1}\|\right\}\leq \frac{8}{\mu}\|g_{t+1}\|.
\end{align*}
Since  $\{x_t\}$ converges to $x^*$ (see Theorem~\ref{thm:3.3}), we conclude that \eqref{ineq:xi-xt<gt+1} holds with $C_1:=8/\mu.$
 

Now, in view of the last inequality in \eqref{eq:subprob-Alg1} and the definition of $\sigma_t$ in Step~2 of Algorithm~1, we have  $\|\sum_{j}(\lambda_{t+1})_j\nabla M^j_{x_t,\alpha\sigma_{t+1}}(x_{t+1})\|\leq \theta\|\Delta x_{t+1}\|^2$, which, combined with   the definitions of $M^j_{x_t,\alpha\sigma_{t+1}}(\cdot)$ and $\tilde g_t$, yields
\begin{equation}\label{ineq:aux-tilde gt-theta}
\left\|\tilde g_t +\sum_j(\lambda_{t+1})_j\nabla^2 \bar F_j(x_t)\Delta x_{t+1}+\alpha\sigma_{t+1}\|\Delta x_{t+1}\|\Delta x_{t+1}\right\|\leq \theta\|\Delta x_{t+1}\|^2.
\end{equation}
Note that in view of the definition of $E^H_j(\cdot)$ in \eqref{ineq:ApproxGrad-Hess}, we have
$$
\|\nabla^2 \bar  F_j(x_t)\|\leq \|\nabla^2 F_j(x_t)\|+\|E^H_j(x_t)\|,
$$
which, combined with  \eqref{ineq:ApproxGrad-Hess12} and the boundedness of $\{x_t\}$, implies that  $\{\nabla^2\bar F_j(x_t)\}$ is also bounded for any $j \in \J$. Hence, there exists $\Lambda>0$ such that $\|\nabla^2\bar F_j(x_t)\|\leq \Lambda$ for any $t\geq 0$ and $j \in \J$. Thus, \eqref{ineq:aux-tilde gt-theta} combined with the triangle inequality and the facts that $\sigma_t\leq \sigma_{max}$ and $\sum_j(\lambda_t)_j=1$, imply that
 \begin{align*}
\left\|\tilde g_t\right\| &\leq   \left\|\sum_j(\lambda_{t+1})_j\nabla^2 \bar F_j(x_t)\Delta x_{t+1}\right\|+\alpha\sigma_{t+1}\|\Delta x_{t+1}\|^2+\theta\|\Delta x_{t+1}\|^2 \\
&\leq \left(\Lambda+(\alpha\sigma_{max}+\theta)\|\Delta x_{t+1}\|\right)\|\Delta x_{t+1}\|.
 \end{align*}
 Since $\{\|\Delta x_{t}\|\}$ converges to zero, it is bounded, and then the above inequality implies that  there exists $\tilde C_2>0$ such that 
 \begin{equation}\label{aux-ineq123}
 \|\tilde g_t\|\leq \tilde C_2\|\Delta x_{t+1}\|,\qquad \forall t\geq 1.
 \end{equation}
Now,  it follows from Lemma~\ref{lem3:basic-gt+1<Deltax} with $(x,,x^+,\lambda^+,\sigma)=(x_t,x_{t+1},\lambda_{t+1},\alpha^{i_t}\sigma_t)$, assumption {\bf (A4)}, and the definitions of $g_t, \tilde g_t, g_{t+1}$, and $\Delta x_{t+1}$ given in \eqref{def:gt-Deltaxt}, that 
 \begin{align*}
\left\|g_{t+1}\right\|&\leq \frac{L+\alpha^{i_t}\sigma_t+2\theta}{2}\|\Delta x_{t+1}\|^2+ \|\Delta x_{t+1}\|\frac{\bar \kappa_{H}}{\alpha^{i_t-1}}\min\left\{ \|\Delta x_{t}\|, \zeta\left\|g_t\right\| \right\}\\
&\leq \frac{8(L+\alpha\sigma_{max}+2\theta)}{\mu^2}\|g_t\|\|\tilde g_t\|+ \frac{4\alpha\bar \kappa_{H}}{\mu} \zeta\left\|g_t\right\|\|\tilde g_t\|,
\end{align*} 
where the last inequality is due to  $\alpha^{i_t}\geq 1$,  $\sigma_{t+1}=\alpha^{i_t-1}\sigma_t$,  the second inequality in \eqref{def:sigmaMax} with $t+1$, and \eqref{eq:Deltaxt-lambdanablaFxt} twice, first with $\lambda=\lambda_{t+1}$ and second with $\lambda=\lambda_t$.
Hence, defining $\tilde C_1:=8(L+\alpha \sigma_{max}+2\theta)/\mu^2+4\alpha\bar\kappa_H\zeta/\mu$, we conclude that 
\begin{equation}\label{ineq:aux-gt+1<gtDelta xt}
\|g_{t+1}\|\leq \tilde C_1\|\tilde g_t\|\|g_t\|.
\end{equation}

Hence, \eqref{ineq:xi-xt<gt+1},  \eqref{aux-ineq123}, \eqref{ineq:aux-gt+1<gtDelta xt}, the fact that $\Delta x_{t+1}=x_{t+1}-x_t$, and the Cauchy-Schwarz inequality imply that 

\begin{align}
 \|x_t-x^*\|&\geq \|x_{t+1}-x_t\|-\|x_{t+1}- x^*\| \geq \frac{1}{\tilde C_2} \left\|\tilde g_t\right\| - C_1 \left\|g_{t+1}\right\| \\
 &= \left(\frac{1}{\tilde C_2}  -C_1\tilde C_1 \left\| g_t\right\|\right)\left\|\tilde g_t\right\|
\end{align}
Since, in view of Corollary~\ref{cor:conv}, $\{g_t\}$ converges to zero,  the above inequality implies that there exists $\hat C_2>0$ such that for any $t$ sufficiently large,
$$
\|x_t- x^*\|\geq \hat C_2\left\|\tilde g_t\right\|.$$
Hence, the above inequality combined with \eqref{aux-ineq123} imply that  \eqref{aux-eq: tilde gt} holds with  $C_2=\max\{\tilde C_2, 1/\hat C_2\}$.
Therefore, in view of \eqref{ineq:xi-xt<gt+1}, \eqref{aux-eq: tilde gt}, and \eqref{ineq:aux-gt+1<gtDelta xt}, we conclude that for all $t$ sufficiently large, we have
$$
\|x_{t+1}-x^*\|\leq C_1\|g_{t+1}\|\leq C_1\tilde C_1\left\|\tilde g_t\right\|\|g_t\|\leq {C_1\tilde C_1C_2}\|g_t\|\|x_t-x^*\|,
$$
which proves \eqref{ineq:xt+1-xbar leq gt xt-xbar} by defining $C_3= C_1\tilde C_1 C_2$.
The proof of the last statement of the theorem follows immediately from \eqref{ineq:xt+1-xbar leq gt xt-xbar} and the fact that $\{g_t\}$ converges to zero, in view of Corollary~\ref{cor:conv}.
\end{proof}

\end{document}
